% Generated by make_pi3b_pi1.py; numeric sources in pi3b_pi1_ssd_sources.json.
\subsection{Physical Raspberry Pi: single thread with an idle USB SSD attached}
\label{app:pi3b-pi1}
A separately preregistered run measures the same twenty frozen pairs and YOLOv8n-seg on a single pinned CPU core of the same board, at the same declared 600~MHz undervoltage. Unlike the four-thread run, a separately powered USB SSD (ADATA SP900, unmounted) is attached throughout; this is a \emph{distinct declared condition} and is not pooled with the four-thread result. The run records the SSD block-I/O counters at every environment snapshot: the written-sector count stays at zero on both USB disks for the whole session, so the disk is verified idle during timing. The session passes every quality gate, including this idle-disk gate, with swap held off, the clock locked at 600~MHz, zero flagged references and an A/A noise floor corresponding to an MDE of 1.002. On the single core the attacked-to-clean total-latency ratio is 1.967 (95\% CI [1.902, 2.030]; all 20 images positive), larger than the four-thread ratio because mask construction benefits from multiple threads while the attacked mask work dominates on one core. A pre-NMS cap reduces the attacked total to a factor of 0.699 (95\% CI [0.682, 0.718]). Median total latency is 10.05~s clean, 20.30~s attacked and 13.77~s capped; the forward pass is 99\% of the clean total. These single-core numbers describe one undervolted board with a declared idle USB device and do not transfer to a 1.2~GHz operating point, to a no-USB condition, or to denial of service.
